% GENERATED FILE. Edit canonical lessons, then run npm run generate:books.
% canonical-source-hash: fnv1a64:8ce38c82319c2665
% canonical-lessons: KA-C356-janapada, KA-C356-samanta, KA-C356-samara, KA-C356-lavani, KA-C356-niti

\chapter{Ballads and Vassal Kings}
\label{ch:ka-ballads-and-vassal-kings}
\hlchaptermodality{356}

\begin{chapteropening}
\textbf{By the end of this chapter:} \emph{I can say folklore, a vassal king, a battle, a ballad and a moral.}

Complete the last lesson of chapter 356: I can say folklore, a vassal king, a battle, a ballad and a moral.
\end{chapteropening}

\section[\texorpdfstring{\kn{ಜಾನಪದ} (jānapada)}{ಜಾನಪದ (jānapada)}]{\kn{ಜಾನಪದ} (jānapada) — folklore; folk}

\label{lesson:KA-C356-janapada}



\begin{quote}
Before the new one: say the Kannada for an inscription, then the Kannada for a manuscript.
\end{quote}

\subsection*{You'll want to know: \kn{ಜಾನಪದ}}
\textbf{\kn{ಜಾನಪದ}} — \emph{jānapada} — \textquotedblleft{}folklore; folk\textquotedblright{}.

\begin{quote}
\kn{ಅಜ್ಜಿಗೆ} \kn{ತುಂಬಾ} \kn{ಜಾನಪದ} \kn{ಹಾಡುಗಳು} \kn{ಗೊತ್ತು}.
\end{quote}

\begin{quote}
\emph{ajjige tumbā jānapada hāḍugaḷu gottu.}
\end{quote}

\begin{quote}
Grandmother knows many folk songs.
\end{quote}

Before a noun it means \textquotedblleft{}folk\textquotedblright{}: \kn{ಜಾನಪದ} \kn{ಕಥೆ}, a folk tale.

\begin{cousinweb}
\textbf{\kn{ಜಾನಪದ}} (\emph{jānapada}) is borrowed from Sanskrit \emph{jānapada}, \textquotedblleft{}living in the country\textquotedblright{}, and as a noun \textquotedblleft{}a country person\textquotedblright{} (Monier-Williams), from \emph{janapada}, \textquotedblleft{}a country, a people\textquotedblright{}.
\end{cousinweb}

\begin{grammarlens}[title={What you've built}]
One more word for telling what happened.
\end{grammarlens}

\subsection*{Your turn}
\begin{itemize}
\raggedright
  \item \emph{Say it:} \emph{jānapada}
  \item \emph{Say it:} \emph{jānapada}, once more
  \item \emph{Say it:} \emph{ajjige tumbā jānapada hāḍugaḷu gottu.} — Grandmother knows many folk songs.
  \item \emph{Recall:} say the Kannada for an inscription, then the Kannada for a manuscript, then say \emph{jānapada} again
\end{itemize}

\begin{tcolorbox}[breakable,colback=teal!4,colframe=teal!35!black,title={Before you move on}]
What does \kn{ಜಾನಪದ} mean? (\textquotedblleft{}Folklore; folk\textquotedblright{}.) Say it once more.
\end{tcolorbox}
\section[\texorpdfstring{\kn{ಸಾಮಂತ} (sāmanta)}{ಸಾಮಂತ (sāmanta)}]{\kn{ಸಾಮಂತ} (sāmanta) — a vassal king}

\label{lesson:KA-C356-samanta}



\begin{quote}
Before the new one: say the Kannada for a manuscript, then the Kannada for folklore.
\end{quote}

\subsection*{You'll want to know: \kn{ಸಾಮಂತ}}
\textbf{\kn{ಸಾಮಂತ}} — \emph{sāmanta} — \textquotedblleft{}a vassal king\textquotedblright{}.

\begin{quote}
\kn{ಆ} \kn{ಸಾಮಂತ} \kn{ಚಕ್ರವರ್ತಿಗೆ} \kn{ಪ್ರತಿ} \kn{ವರ್ಷ} \kn{ಕಪ್ಪ} \kn{ಕೊಡುತ್ತಿದ್ದನು}.
\end{quote}

\begin{quote}
\emph{ā sāmanta cakravartige prati varṣa kappa koḍuttiddanu.}
\end{quote}

\begin{quote}
That vassal king paid tribute to the emperor every year.
\end{quote}

Plural: \kn{ಸಾಮಂತರು}, \emph{sāmantaru}.

A local ruler under a greater king.

\begin{cousinweb}
\textbf{\kn{ಸಾಮಂತ}} (\emph{sāmanta}) is borrowed from Sanskrit \emph{sāmanta}, \textquotedblleft{}being on all sides\textquotedblright{}, and as a noun \textquotedblleft{}a neighbour\textquotedblright{}, later \textquotedblleft{}a vassal\textquotedblright{} (CDIAL 13348).
\end{cousinweb}

\begin{grammarlens}[title={What you've built}]
One more word for telling what happened.
\end{grammarlens}

\subsection*{Your turn}
\begin{itemize}
\raggedright
  \item \emph{Say it:} \emph{sāmanta}
  \item \emph{Say it:} \emph{sāmanta}, once more
  \item \emph{Say it:} \emph{ā sāmanta cakravartige prati varṣa kappa koḍuttiddanu.} — That vassal king paid tribute to the emperor every year.
  \item \emph{Recall:} say the Kannada for a manuscript, then the Kannada for folklore, then say \emph{sāmanta} again
\end{itemize}

\begin{tcolorbox}[breakable,colback=teal!4,colframe=teal!35!black,title={Before you move on}]
What does \kn{ಸಾಮಂತ} mean? (\textquotedblleft{}A vassal king\textquotedblright{}.) Say it once more.
\end{tcolorbox}
\section[\texorpdfstring{\kn{ಸಮರ} (samara)}{ಸಮರ (samara)}]{\kn{ಸಮರ} (samara) — a battle; combat}

\label{lesson:KA-C356-samara}



\begin{quote}
Before the new one: say the Kannada for folklore, then the Kannada for a vassal king.
\end{quote}

\subsection*{You'll want to know: \kn{ಸಮರ}}
\textbf{\kn{ಸಮರ}} — \emph{samara} — \textquotedblleft{}a battle; combat\textquotedblright{}.

\begin{quote}
\kn{ಆ} \kn{ಸಮರದಲ್ಲಿ} \kn{ಎರಡು} \kn{ಸಾಮ್ರಾಜ್ಯಗಳು} \kn{ಕಾದಾಡಿದವು}.
\end{quote}

\begin{quote}
\emph{ā samaradalli eraḍu sāmrājyagaḷu kādāḍidavu.}
\end{quote}

\begin{quote}
In that battle two empires fought each other.
\end{quote}

Plural: \kn{ಸಮರಗಳು}, \emph{samaragaḷu}.

A bookish word; \kn{ಸಮರ} \kn{ಸಾರು}, to declare war.

\begin{cousinweb}
\textbf{\kn{ಸಮರ}} (\emph{samara}) is borrowed from Sanskrit \emph{samara}, \textquotedblleft{}coming together\textquotedblright{}, and so \textquotedblleft{}a hostile encounter, a battle\textquotedblright{} (Monier-Williams).
\end{cousinweb}

\begin{grammarlens}[title={What you've built}]
One more word for telling what happened.
\end{grammarlens}

\subsection*{Your turn}
\begin{itemize}
\raggedright
  \item \emph{Say it:} \emph{samara}
  \item \emph{Say it:} \emph{samara}, once more
  \item \emph{Say it:} \emph{ā samaradalli eraḍu sāmrājyagaḷu kādāḍidavu.} — In that battle two empires fought each other.
  \item \emph{Recall:} say the Kannada for folklore, then the Kannada for a vassal king, then say \emph{samara} again
\end{itemize}

\begin{tcolorbox}[breakable,colback=teal!4,colframe=teal!35!black,title={Before you move on}]
What does \kn{ಸಮರ} mean? (\textquotedblleft{}A battle; combat\textquotedblright{}.) Say it once more.
\end{tcolorbox}
\section[\texorpdfstring{\kn{ಲಾವಣಿ} (lāvaṇi)}{ಲಾವಣಿ (lāvaṇi)}]{\kn{ಲಾವಣಿ} (lāvaṇi) — a ballad}

\label{lesson:KA-C356-lavani}



\begin{quote}
Before the new one: say the Kannada for a vassal king, then the Kannada for a battle.
\end{quote}

\subsection*{You'll want to know: \kn{ಲಾವಣಿ}}
\textbf{\kn{ಲಾವಣಿ}} — \emph{lāvaṇi} — \textquotedblleft{}a ballad\textquotedblright{}.

\begin{quote}
\kn{ಹಳ್ಳಿಯ} \kn{ಕಲಾವಿದ} \kn{ಒಂದು} \kn{ಹಳೆಯ} \kn{ಲಾವಣಿ} \kn{ಹಾಡಿದನು}.
\end{quote}

\begin{quote}
\emph{haḷḷiya kalāvida ondu haḷeya lāvaṇi hāḍidanu.}
\end{quote}

\begin{quote}
The village artist sang an old ballad.
\end{quote}

Plural: \kn{ಲಾವಣಿಗಳು}, \emph{lāvaṇigaḷu}.

A sung story of a hero, a battle or a famous event.

\begin{cousinweb}
\textbf{\kn{ಲಾವಣಿ}} (\emph{lāvaṇi}) is usually traced to Marathi \emph{lāvaṇī}, a kind of folk song; how the word came into Kannada is not settled.
\end{cousinweb}

\begin{grammarlens}[title={What you've built}]
One more word for telling what happened.
\end{grammarlens}

\subsection*{Your turn}
\begin{itemize}
\raggedright
  \item \emph{Say it:} \emph{lāvaṇi}
  \item \emph{Say it:} \emph{lāvaṇi}, once more
  \item \emph{Say it:} \emph{haḷḷiya kalāvida ondu haḷeya lāvaṇi hāḍidanu.} — The village artist sang an old ballad.
  \item \emph{Recall:} say the Kannada for a vassal king, then the Kannada for a battle, then say \emph{lāvaṇi} again
\end{itemize}

\begin{tcolorbox}[breakable,colback=teal!4,colframe=teal!35!black,title={Before you move on}]
What does \kn{ಲಾವಣಿ} mean? (\textquotedblleft{}A ballad\textquotedblright{}.) Say it once more.
\end{tcolorbox}
\section[\texorpdfstring{\kn{ನೀತಿ} (nīti)}{ನೀತಿ (nīti)}]{\kn{ನೀತಿ} (nīti) — a moral; a principle}

\label{lesson:KA-C356-niti}



\begin{quote}
Before the new one: say the Kannada for a battle, then the Kannada for a ballad.
\end{quote}

\subsection*{You'll want to know: \kn{ನೀತಿ}}
\textbf{\kn{ನೀತಿ}} — \emph{nīti} — \textquotedblleft{}a moral; a principle\textquotedblright{}.

\begin{quote}
\kn{ಈ} \kn{ಕಥೆಯ} \kn{ನೀತಿ} \kn{ಏನು}?
\end{quote}

\begin{quote}
\emph{ī katheya nīti ēnu?}
\end{quote}

\begin{quote}
What is the moral of this story?
\end{quote}

\kn{ನೀತಿ} \kn{ಕಥೆ} is a fable. \kn{ನೀತಿ} also means a policy: \kn{ಸರ್ಕಾರದ} \kn{ನೀತಿ}.

\begin{cousinweb}
\textbf{\kn{ನೀತಿ}} (\emph{nīti}) is borrowed from Sanskrit \emph{nīti}, \textquotedblleft{}conduct, guidance, wise behaviour\textquotedblright{} (CDIAL 7546), from \emph{nī}, \textquotedblleft{}to lead\textquotedblright{}.
\end{cousinweb}

\begin{grammarlens}[title={What you've built}]
One more word for telling what happened.
\end{grammarlens}

\subsection*{Your turn}
\begin{itemize}
\raggedright
  \item \emph{Say it:} \emph{nīti}
  \item \emph{Say it:} \emph{nīti}, once more
  \item \emph{Say it:} \emph{ī katheya nīti ēnu?} — What is the moral of this story?
  \item \emph{Recall:} say the Kannada for a battle, then the Kannada for a ballad, then say \emph{nīti} again
\end{itemize}

\begin{tcolorbox}[breakable,colback=teal!4,colframe=teal!35!black,title={Before you move on}]
What does \kn{ನೀತಿ} mean? (\textquotedblleft{}A moral; a principle\textquotedblright{}.) Say it once more.
\end{tcolorbox}
